
% \section{Performance Metrics}
% In evaluating the performance of the earlier mentioned scenarios, the following metrics were used :-

% \subsection{Precision}
% This metric brings into perception, the number of correctly made positive predictions. It is the ratio that exist between \ac{TP} and all positives, for our given problem, precision can be defined as the sum of \ac{TP} across all classes divided by the sum of \ac{TP} and \ac{FP} across all classes. This can be expressed mathematically as:-

% \begin{equation}
%     Precision = \frac{TP}{TP + FP}
% \end{equation}

% \subsection{Recall}

% \subsection{Accuracy}




% \subsection{F1-score}




merging\documentclass[11pt]{report}
\usepackage[utf8]{inputenc}
\usepackage[a4paper,width=190mm, height = 260mm, left=25mm,right = 25mm, top=25mm,bottom=25mm,bindingoffset=10mm]{geometry}

% \usepackage{geometry}
%  \geometry{
%  a4paper,
%  total={170mm,257mm},
%  left=20mm,
%  top=20mm,
%  }
\usepackage{cite}
\usepackage{amsmath,amssymb,amsfonts}
\usepackage{graphicx}
\usepackage{textcomp}
\usepackage{color, colortbl}
\definecolor{Gray}{gray}{0.9}
\usepackage{multirow}
\usepackage[usenames,dvipsnames]{xcolor}
\usepackage{longtable}
\usepackage{booktabs}
\usepackage{tabularx}
\usepackage{amsmath}
\interdisplaylinepenalty=2500
\usepackage[cmintegrals]{newtxmath}
\usepackage[caption=false]{subfig}
\usepackage{filecontents}
\usepackage{float}
\usepackage{acronym}
\usepackage{CJKutf8}
\usepackage[caption=false]{subfig}
\usepackage[titles]{tocloft} 

\newlength\mylength
\renewcommand\cftchappresnum{\chaptername~}

\renewcommand{\bibname}{References} 
\renewcommand\cftchapaftersnum{:~}
\settowidth\mylength{\cftchappresnum\cftchapaftersnum\quad}
\addtolength\cftchapnumwidth{\mylength}